\section{Neural Style Transfer (NST)}

The Python part of this exercise is available in \myhref{https://github.com/LosDelDGIIM/LosDelDGIIM.github.io/blob/main/subjects/Erasmus-DUE/GKI/Exercises/Exercise11.ipynb}{this Jupyter Notebook}.\\

\begin{ejercicio}[Basic Concepts of Training]
    Explain briefly and clearly:
    \begin{itemize}
        \item What is a train-test split?
        
        A train-test split is a technique used in machine learning to divide a dataset into two separate subsets: one for training the model and another for testing its performance. The training set is used to fit the model, while the test set is used to evaluate how well the model generalizes to unseen data. This helps prevent overfitting and provides an estimate of the model's accuracy on new data.
        \item What is an epoch?
        
        An epoch is a complete pass through the entire training dataset during the training process of a machine learning model. In other words, one epoch means that every sample in the training dataset has been used once to update the model's parameters. Multiple epochs are often required to achieve good performance, as the model learns and improves with each pass through the data.
        \item What is a batch?
        
        A batch is a subset of the training dataset that is processed together in a single forward and backward pass through the neural network. Batching allows for more efficient computation and can lead to more stable gradient estimates during training.
        \item What is loss?
        
        Loss is a measure of how well the model's predictions match the actual target values. It quantifies the error between the predicted and true outputs, guiding the optimization process to minimize this error.
        \item What is a loss-epoch curve?
        
        A loss-epoch curve is a graphical representation showing how the loss of the model changes over multiple epochs of training. It helps visualize the training progress and can indicate whether the model is learning effectively or if issues like overfitting or underfitting are present.
        \item What is the purpose of the learning rate, and how does it affect the loss-epoch curve?
        
        The learning rate determines how much to adjust the model's parameters in response to the estimated error during training. A high learning rate can cause rapid convergence but may overshoot optimal solutions, while a low learning rate ensures more stable convergence but can result in slower training. The learning rate directly influences the shape and behavior of the loss-epoch curve.
    \end{itemize}
    
\end{ejercicio}


\begin{ejercicio}[RNN]
    How does an RNN work with data that has a sequence (e.g., sentences or time series)? What types of feedback are there? Why do older RNNs suffer from the problem of signals disappearing or exploding during training? And how do LSTMs and GRUs solve this problem?

    Recurrent Neural Networks (RNNs) are designed to handle sequential data by maintaining a hidden state that captures information from previous time steps. This allows RNNs to process sequences of varying lengths, such as sentences or time series, by passing the hidden state from one time step to the next. There are three main types of feedback in RNNs:
    \begin{itemize}
        \item \ul{Direct feedback}: The output of a neuron is fed back into neurons of the next layer.
        \item \ul{Lateral feedback}: The output of a neuron is fed back into neurons of the same layer.
        \item \ul{Indirect feedback}: The output of a neuron is fed back into neurons of previous layers.
    \end{itemize}

    Older RNNs suffer from the problem of vanishing or exploding gradients during training, which occurs when the gradients used to update the model's parameters either become too small (vanish) or too large (explode). This makes it difficult for the model to learn long-term dependencies in the data. Long Short-Term Memory (LSTM) networks and Gated Recurrent Units (GRUs) address this issue by introducing gating mechanisms that regulate the flow of information through the network, allowing them to maintain and update their hidden states more effectively over longer sequences.

    % // TODO: Corregir
\end{ejercicio}

\begin{ejercicio}[GAN]
    How do the generator and discriminator compete against each other in a GAN? And how does this ultimately result in a realistic image?

    In a Generative Adversarial Network (GAN), the generator and discriminator are two neural networks that compete against each other in a zero-sum game. The generator's goal is to create realistic data samples (e.g., images) that can fool the discriminator into thinking they are real, while the discriminator's goal is to correctly distinguish between real and generated samples. During training, the generator improves its ability to produce realistic samples by learning from the feedback provided by the discriminator, while the discriminator improves its ability to detect fake samples. This adversarial process continues until the generator produces samples that are indistinguishable from real data, resulting in highly realistic images.
\end{ejercicio}

\begin{ejercicio}[Autoencoder]
    How is an autoencoder structured? What kinds of problems can it solve?

    An autoencoder is a type of neural network designed to learn efficient representations of data, typically for the purpose of dimensionality reduction or feature learning. It consists of two main components: the encoder and the decoder. The encoder compresses the input data into a lower-dimensional latent space representation, while the decoder reconstructs the original data from this compressed representation. Autoencoders can be used for various tasks, including denoising, anomaly detection, and data compression.
\end{ejercicio}

\begin{ejercicio}[Latent Space]
    What is the latent space of an autoencoder?

    The latent space of an autoencoder refers to the lower-dimensional representation of the input data that is learned by the encoder. It captures the essential features and structure of the data while discarding redundant information. The latent space can be used for various purposes, such as visualization, clustering, or as a compact representation for downstream tasks. It allows for efficient storage and manipulation of data while preserving its underlying characteristics.
\end{ejercicio}